%------------------------------------------
\chapter{Revisão da literatura}\label{cap:desenvolvimento}
\label{cha:tecboot}
% -----------------------------------------
Esta seção apresenta a Regressão Linear Múltipla (RLM) e a técnica \textit{bootstrap}. Adicionalmente, são abordados conceitos como a teoria da probabilidade, distribuição empírica, princípio \textit{plug-in} e técnicas clássicas de reamostragem, fundamentais para a compreensão.
\begin{comment}
%-------------------------------------------------------
\section{Regressão Linear Múltipla}
%-------------------------------------------------------
A RLM é empregada quando é necessário estabelecer uma relação envolvendo mais de uma variável independente (\textit{\gls{x2}})
\begin{equation}
    \gls{y} = x_1b_1 + x_2b_2 +\ldots+ x_nb_n
\end{equation}
\begin{equation}
    b_j \in \gls{R}, \quad j = 1,2,\ldots,n
\end{equation}
ou, na forma matricial,
\begin{equation}
    y = [x_1 \quad x_2 \quad \ldots \quad \gls{x_n}]\left[\begin{array}{c}
b_1\\b_2\\\vdots \\b_n
\end{array}\right] \triangleq \textbf{xb}
\end{equation}

O procedimento de estimação é denominado regressão linear múltipla~\cite{KRAMER:1998,MARTENS:1989,Filzmoser2021,Schielzeth2020,revelas2024,paleologo2020,Dar2024,rathnayake2022}. Dispondo de dados de calibração
\begin{equation}
    \{(x^i,\gls{y^i}), \quad i = 1,2,\ldots,M \}
\end{equation}
ou, na forma matricial
\begin{equation}
X_{cal} =
\left[\begin{array}{c}
x^1\\x^2\\\vdots \\x^M
\end{array}\right]
\quad y_{cal} =
\left[\begin{array}{c}
y^1\\y^2\\\vdots \\y^M
\end{array}\right]
\end{equation}
onde a melhor estimativa no sentido de mínimos quadrados para o vetor de coeficientes \textbf{b} é dada pela pseudo-inversa de Moore-Penrose~\cite{DRAPER:1998,TIMM:2002}:
\begin{equation}
    \widehat{\textbf{b}} = \gls{X_cal^+} \textbf{y}_{cal}\triangleq (\textbf{X}_{cal}^{\textbf{T}} \textbf{X}_{cal})^{-1} \textbf{X}_{cal}^{\textbf{T}} \textbf{y}_{cal}
\end{equation}

Segundo Alin e Serneels~\cite{Alin2017} e Filzmoser e Nordhausen~\cite{Filzmoser2013}, métodos robustos se tornam necessários quando há violação das premissas da regressão linear múltipla.


As premissas de normalidade dos resíduos devem ser verificadas em modelos de calibração. Testes robustos de normalidade, especialmente apropriados para dados quimiométricos - que frequentemente apresentam estruturas complexas - são discutidos em estudos anteriores~\cite{Alin2017,Filzmoser2013}.
\end{comment}


\section{Regressão Linear Múltipla}

A RLM é uma técnica estatística empregada para modelar a relação linear entre uma variável dependente (resposta) ($\gls{y}$) e duas ou mais variáveis independentes (preditoras) ($\mathbf{X}$). O modelo, que inclui um termo de intercepto ($\beta_0$) e um termo de erro ($\varepsilon$), é formalizado como:
\begin{equation}
    \label{eq:rlm-modelo}
    \gls{y} = \beta_0 + x_1\beta_1 + x_2\beta_2 +\ldots+ x_n\beta_n + \varepsilon
\end{equation}
onde $\beta_j \in \mathbb{R}$ representa o coeficiente angular de cada preditor $x_j$, para $j = 1,2,\ldots,n$.

Dispondo de um conjunto de $M$ dados de calibração $\{(x^i,\gls{y^i}), \quad i = 1,2,\ldots,M \}$, o modelo pode ser expresso na forma matricial:
\begin{equation}
    \mathbf{y} = \mathbf{X}\mathbf{\beta} + \mathbf{\varepsilon}
\end{equation}
onde o vetor de resposta $\mathbf{y}$ e a matriz de calibração $\mathbf{X}$ (incluindo a coluna de uns para o intercepto) são definidos por:
\begin{equation}
    \mathbf{X} =
    \left[\begin{array}{ccccc}
    1 & x_{1}^1 & x_{2}^1 & \ldots & x_{n}^1 \\
    1 & x_{1}^2 & x_{2}^2 & \ldots & x_{n}^2 \\
    \vdots & \vdots & \vdots & \ddots & \vdots \\
    1 & x_{1}^M & x_{2}^M & \ldots & x_{n}^M
    \end{array}\right]
    \quad \mathbf{y} =
    \left[\begin{array}{c}
    y^1\\y^2\\\vdots \\y^M
    \end{array}\right]
\end{equation}
O vetor de coeficientes populacionais é $\mathbf{\beta} = [\beta_0, \beta_1, \ldots, \beta_n]^{\textbf{T}}$.

\subsection{Estimação dos Coeficientes}
O procedimento de estimação é conhecido como \gls{MQO}. O MQO busca o vetor de coeficientes estimados $\widehat{\mathbf{\beta}}$ que minimiza a \gls{SQR}.

A melhor estimativa $\widehat{\mathbf{\beta}}$ no sentido de mínimos quadrados é dada pela solução da pseudo-inversa de Moore-Penrose \cite{DRAPER:1998,TIMM:2002}:
\begin{equation}
    \label{eq:rlm-pseudo}
    \widehat{\mathbf{\beta}} = \mathbf{X}^+ \mathbf{y} \triangleq (\mathbf{X}^{\textbf{T}} \mathbf{X})^{-1} \mathbf{X}^{\textbf{T}} \mathbf{y}
\end{equation}
onde $\mathbf{X}^{+}$ é a pseudo-inversa da matriz $\mathbf{X}$.


\subsection{Premissas e Métodos Robustos}
A validade das estimativas obtidas pelo MQO depende do cumprimento das premissas clássicas da regressão linear. Em dados espectrais, a colinearidade entre variáveis e a presença de outliers violam com frequência esses pressupostos~\cite{Alin2017,Filzmoser2013}.

Métodos robustos tornam-se indispensáveis quando há violação das premissas da regressão linear múltipla, especialmente na presença de \textit{outliers} \cite{Alin2017,Filzmoser2013}. A verificação da normalidade dos resíduos, em particular, é crucial e discutida em detalhes em estudos anteriores \cite{Alin2017,Filzmoser2013}.


% ----------------------------------------------------
\section{A Técnica de \textit{Bootstrap}}
% ----------------------------------------------------
%A inferência estatística estabelece propriedades para a população \textbf{X} a partir da análise de uma amostra aleatória dessa população~\cite{DRAPER:1998,JOHNSON:2001,CHATTERJEE:2006}. O método \textit{bootstrap} é um procedimento computacional que fornece a medida do nível de precisão da inferência estatística, como por exemplo, média, desvio padrão e erro padrão~\cite{yeon2023,alfons2022,wan2023,du2024,jensen2024}. A ideia básica do \textit{bootstrap} é inferir estatisticamente as propriedades de interesse a partir da análise de amostras aleatórias extraídas da população original sem a necessidade de conhecer a distribuição de dados a priori, o que constitui significativa vantagem sobre outras técnicas~\cite{HARDLE:2007,TIMM:2002}.

A inferência estatística estabelece propriedades para a população \textbf{X} a partir da análise
de uma amostra aleatória dessa população~\cite{DRAPER:1998,JOHNSON:2001,CHATTERJEE:2006}. O método \textit{bootstrap} é um procedimento computacional que fornece a medida do nível de precisão da inferência estatística, como por exemplo, média, desvio padrão e erro padrão~\cite{yeon2023,alfons2022,wan2023,du2024,jensen2024}.

A ideia básica do \textit{bootstrap} é inferir estatisticamente as propriedades de interesse a partir da análise de amostras aleatórias extraídas da população original sem a necessidade de conhecer a distribuição de dados a priori, o que dispensa a suposição prévia de uma distribuição paramétrica dos dados~\cite{HARDLE:2007,TIMM:2002}.

Quando o parâmetro populacional de interesse é desconhecido, procura-se estimá-lo por
reamostragem aleatória. Inicialmente, considera-se a \textbf{população} 
$\mathbf{x} = [x_1, x_2, \ldots, x_N]$. A partir dela, seleciona-se uma 
\textbf{amostra} $\mathbf{X} = [x_1, x_2, \ldots, x_n]$, formada por $n$ elementos 
independentes retirados da população. 

%Para avaliar a variabilidade das estimativas, aplica-se o método de \textbf{\textit{bootstrap}}, que consiste em gerar $m$ \textbf{reamostras} $\mathbf{X}^{*1}, \mathbf{X}^{*2}, \ldots, \mathbf{X}^{*m}$, obtidas por sorteio com reposição a partir da amostra original. Formalmente, cada reamostra é definida como \[ \mathbf{X}^{*i} \sim \text{\textit{Bootstrap}}(\mathbf{X}), \quad i = 1,2,\ldots,m,\] onde $\text{\textit{Bootstrap}}(\mathbf{X})$ denota a distribuição empírica gerada pela amostra. Em seguida, calcula-se uma estatística $s(\mathbf{X}^{*i})$ para cada reamostra, gerando uma distribuição empírica das estimativas. Esse processo é ilustrado na Figura~\ref{fig:bootstrap}.

Para avaliar a variabilidade das estimativas, aplica-se o método de \textbf{\textit{bootstrap}}, que consiste em gerar $m$ \textbf{reamostras} 
$\mathbf{X}^{*1}, \mathbf{X}^{*2}, \ldots, \mathbf{X}^{*m}$, obtidas por sorteio com reposição a partir da amostra original. Formalmente, cada reamostra é definida como
\[
\mathbf{X}^{*i} \sim \text{\textit{Bootstrap}}(\mathbf{X}), \quad i = 1,2,\ldots,m,
\]
onde $\text{\textit{Bootstrap}}(\mathbf{X})$ denota a distribuição empírica gerada pela amostra.

Em seguida, calcula-se uma estatística $s(\mathbf{X}^{*i})$ para cada reamostra, 
gerando uma distribuição empírica das estimativas. Esse processo é ilustrado 
na Figura~\ref{fig:bootstrap}.

\begin{figure}[H]
    \centering
		\caption{Exemplo do método \textit{bootstrap}.}
		\resizebox{0.70\textwidth}{!}{
    \begin{tikzpicture}[
        box/.style={rectangle, draw, minimum width=2.5cm, minimum height=1cm, align=center},
        arrow/.style={-Stealth, thick}
    ]
    
    % Nodo População
    \node[box] (Pop) at (0,2) {$\mathbf{x} = [x_1, x_2, \dots, x_N]$ \\ População};
    
    % Nodo Amostra
    \node[box] (X) at (0,0) {$\mathbf{X} = [x_1, x_2, \dots, x_n]$ \\ Amostra};
    
    % Nodos de Reamostras
    \node[box] (X1) at (-3,-2) {$\mathbf{X}^{*1}$};
    \node[box] (X2) at (0,-2) {$\mathbf{X}^{*2}$};
    \node[box, dotted] (dots) at (3,-2) {$\dots$};
    \node[box] (Xm) at (6,-2) {$\mathbf{X}^{*m}$};
    
    % Nodos de Estatísticas
    \node[box] (sX1) at (-3,-4) {$s(\mathbf{X}^{*1})$};
    \node[box] (sX2) at (0,-4) {$s(\mathbf{X}^{*2})$};
    \node[box, dotted] (sdots) at (3,-4) {$\dots$};
    \node[box] (sXm) at (6,-4) {$s(\mathbf{X}^{*m})$};
    
    % Setas População -> Amostra
    \draw[arrow] (Pop) -- (X) node[midway, right] {Amostragem};
    
    % Setas Amostra -> Reamostras
    \draw[arrow] (X) -- (X1) node[midway, above left] {Reamostras};
    \draw[arrow] (X) -- (X2);
    \draw[arrow] (X) -- (dots);
    \draw[arrow] (X) -- (Xm) node[midway, above right] {Reamostras};
    
    % Setas Reamostras -> Estatísticas
    \draw[arrow] (X1) -- (sX1) node[midway, left] {Estatísticas};
    \draw[arrow] (X2) -- (sX2);
    \draw[arrow] (dots) -- (sdots);
    \draw[arrow] (Xm) -- (sXm) node[midway, right] {Estatísticas};
    
    \end{tikzpicture}
		}
\label{fig:bootstrap}
\legend{\footnotesize Fonte:~\oautor}
\end{figure}

%A Figura~\ref{fig:bootstrap} ilustra o procedimento de \textit{bootstrap} em quatro etapas. O primeiro passo (primeiro retângulo) representa a população completa, denotada por $\mathbf{x} = [x_1, x_2, \dots, x_N]$, da qual é extraída uma amostra através do processo de amostragem. O segundo passo (segundo retângulo) corresponde ao conjunto de dados observados, denotado por $\mathbf{X} = [x_1, x_2, \dots, x_n]$. O terceiro passo (terceiro conjunto de retângulos) corresponde à extração de amostras \textit{bootstrap}, obtidas por reamostragem com reposição a partir de $\mathbf{X}$. O quarto passo (quarto conjunto de retângulos) mostra as replicações \textit{bootstrap}, nas quais uma função estatística $s(\gls{gamma})$ é recalculada $m$ vezes, sendo $m$ o número total de replicações.

A Figura~\ref{fig:bootstrap} ilustra o procedimento de \textit{bootstrap} em quatro etapas. O primeiro passo (primeiro retângulo) representa a população completa, denotada por $\mathbf{x} = [x_1, x_2, \dots, x_N]$, da qual é extraída uma amostra através do processo de amostragem. O segundo passo (segundo retângulo) corresponde ao conjunto de dados observados, denotado por $\mathbf{X} = [x_1, x_2, \dots, x_n]$.

O terceiro passo (terceiro conjunto de retângulos) corresponde à extração de amostras \textit{bootstrap}, obtidas por reamostragem com reposição a partir de $\mathbf{X}$. O quarto passo (quarto conjunto de retângulos) mostra as replicações \textit{bootstrap}, nas quais uma função estatística $s(\gls{gamma})$ é recalculada $m$ vezes, sendo $m$ o número total de replicações.

Nas seções subsequentes, são discutidos conceitos fundamentais relacionados a amostras aleatórias e probabilísticas, bem como a distribuição empírica e o princípio \textit{plug-in}. A apresentação desses tópicos tem como objetivo fornecer a base teórica necessária para a compreensão dos procedimentos de inferência estatística realizados a partir de amostras \textit{bootstrap}.

%------------------------------------------------
\subsection{\textit{Bootstrap}}
\label{sec:boot}
\index{A Técnica de \textit{Bootstrap}!\textit{Bootstrap}}
%-----------------------------------------------
Uma amostra aleatória $\gls{x_al} = [x_{al1},\ldots, x_{aln}]$ é extraída da distribuição de probabilidade \textbf{D}. O parâmetro de interesse $\theta = \gls{t(d)}$ é desconhecido e será estimado com base na amostra observada $\gls{x_al}$, resultando no estimador $\widehat{\theta} = t(\gls{x_al})$.

%-----------------------------------------------
\subsection{Estimativa \textit{Bootstrap}}
\label{sec:estboot}
\index{A Técnica de \textit{Bootstrap}!Estimativa \textit{Bootstrap}}
%-----------------------------------------------
O método \textit{bootstrap} necessita inicialmente das amostras \textit{bootstrap} para depois aplicar as replicações \textit{bootstrap}, como foi visto na Figura~\ref{fig:bootstrap}. Suponha a distribuição empírica $\widehat{\textbf{d}}$, com probabilidade $\frac{1}{n}$ em cada valor de $x_i$, tendo $i = 1, 2,\ldots, n$. A amostra \textit{bootstrap} é definida como sendo a amostra aleatória de tamanho $n$ extraída de $\widehat{\textbf{d}}$, representada por $\gls{x^*} =[x_1^*, x_2^*,\ldots, x_n^*]$~\cite{EFRON:1994}. As amostras \textit{bootstrap} são obtidas por reamostragem com reposição, podendo haver repetições dos elementos da amostra original $[x_1, x_2,\ldots, x_n]$.

As replicações \textit{bootstrap} são representadas por
\begin{equation}
    \widehat{\theta} = s(\gls{x^*}),
\end{equation}
onde $s(\gls{gamma})$ é a função estatística a ser aplicada sobre as amostras \textit{bootstrap}. Por exemplo, se $s(x)$ for a média amostral $\overline{x}$, então $s(\gls{x^*})$ é a média da amostra \textit{bootstrap}
\begin{equation}
    \overline{x^*} = \frac{1}{n} \sum_{i=1}^{n} x_i^*.
\end{equation}

\medskip
%--------------------------------------------------
\begin{minipage}{0.80\textwidth}
\begin{algorithm2e}[H]
 \dontprintsemicolon %\dontprintsemicolon
 \linesnumbered%\linesnumbered
 %%\nocaptionofalgo
 \SetLine
 \BlankLine
 \Entrada{Amostra original $x = (x_1, x_2, \ldots, x_n)$, estatística de interesse $\theta = s(x)$, número de replicações \textit{bootstrap} $m$}
 \Saida{Estimativa \textit{bootstrap} \gls{hat_theta^*}, erro padrão \textit{bootstrap}, intervalo de confiança \textit{bootstrap}}
 \BlankLine
\Inicio{
  Selecione a quantidade de amostras independentes ~$m$,~\{$x^{*1},x^{*2},\ldots,x^{*m}$\}, tendo tamanho
 $m$ e podendo haver repetições de  $x$.\;
Avaliar a replicação \textit{bootstrap} para amostra \textit{bootstrap}, $\widehat{\theta}^*(b)= s(x^{*b})$ $\quad$ $b = 1,2,\ldots ,m$. \;}
\caption{Estimativa das replicações \textit{bootstrap}. \label{alg:bootstrap}}
\end{algorithm2e}
\legend{\footnotesize Fonte:~\oautor}
\end{minipage}

O Algoritmo~\ref{alg:bootstrap} \cite{EFRON:1993} mostra a descrição do procedimento para
utilizar a estimativa \textit{bootstrap}, replicando o estimador $\widehat{\theta} = s(x)$ a partir dos dados $x$ observados.

%-------------------------------------------------
\subsection{O Número de Replicações \textit{Bootstrap} \textit{m}}
\label{sec:numreplboot}
\index{A Técnica de \textit{Bootstrap}!O Número de Replicações \textit{Bootstrap} $m$}
%-------------------------------------------------
A estimativa \textit{bootstrap} ideal é obtida através de uma boa escolha do número de replicações
$\gls{m} = \infty$, com o $\gls{hat_se}_{\infty}$ igual à estimativa \textit{plug-in} $se_{\widehat{\textbf{d}}}(\widehat{\theta}^*)$~\cite{EFRON:1993}.
O número de replicações \textit{bootstrap} influencia diretamente o tempo computacional
gasto ao processar as replicações, principalmente se a função $s(x)$ do estimador $\widehat{\theta} = s(x)$ for muito
complexa. Portanto, se não for possível obter infinitas replicações, deve-se utilizar o maior número possível de replicações.

%---------------------------------------------------------------------------------------------------------

%É necessário que estimadores apresentem bom comportamento: baixo erro de predição e pequeno desvio padrão. A estimativa \textit{bootstrap} para o erro padrão geralmente apresenta erro de predição relativamente pequeno. A estimativa \textit{bootstrap} ideal $\widehat{se}_{\infty}$ possui baixa variabilidade entre estimativas com erro de predição quase nulo de $se_{\widehat{\textbf{d}}}(\widehat{\theta^*})$, pelo menos em um sentido assintótico. Esta propriedade decorre do fato de que $\widehat{se}_m$ é a estimativa \textit{plug-in} de $se_{\widehat{\textbf{d}}}(\widehat{\theta^*})$. Não é difícil mostrar que $\widehat{se}_m$ sempre possui desvio padrão maior que $\widehat{se}_{\infty}$~\cite{EFRON:1993}. A questão que se coloca é: \textbf{o quanto maior}?

É necessário que estimadores apresentem bom comportamento: baixo erro de predição e pequeno desvio padrão. A estimativa \textit{bootstrap} para o erro padrão geralmente apresenta erro de predição relativamente pequeno. A estimativa \textit{bootstrap} ideal $\widehat{se}_{\infty}$ possui baixa variabilidade entre estimativas com erro de predição quase nulo de $se_{\widehat{\textbf{d}}}(\widehat{\theta^*})$, pelo menos em um sentido assintótico.

Esta propriedade decorre do fato de que $\widehat{se}_m$ é a estimativa \textit{plug-in} de $se_{\widehat{\textbf{d}}}(\widehat{\theta^*})$. Não é difícil mostrar que $\widehat{se}_m$ sempre possui desvio padrão maior que $\widehat{se}_{\infty}$~\cite{EFRON:1993}. A questão que se coloca é: \textbf{o quanto maior}?

%--------------------------------------------------------------------------------------------------

Uma resposta aproximada, porém satisfatória, pode ser encontrada utilizando o coeficiente de variação de $\widehat{se}_{m}$ em relação à sua expectativa. O aumento da variabilidade decorrente da interrupção do processo após $m$ replicações \textit{bootstrap}, em vez de permitir que $m$ tenda ao infinito, resulta em um aumento do coeficiente de variação:
\begin{equation}
    cv(\widehat{se}_{m}) = \sqrt{[cv(\widehat{se}_\infty)]^2 + \frac{E(\Delta) + 2}{4m}}.
\label{eq:2.19}
\end{equation}

O parâmetro \gls{Delta} mede o tamanho da distribuição de $\widehat{\theta}^*$: $\Delta$ é zero para a distribuição normal; na prática, $\Delta$ geralmente é maior que 10. O coeficiente de variação na Equação~(\ref{eq:2.19}) incorpora tanto a variabilidade devida à reamostragem \textit{bootstrap} quanto a variabilidade devida à amostragem da população original.


A estimativa ideal $\widehat{se}_\infty = se_{\widehat{\mathbf{d}}}(\widehat{\theta}^*)$ não é perfeita~\cite{EFRON:1993}. Pode haver variabilidade estatisticamente significativa na estimativa $se_{\mathbf{d}}(\widehat{\theta})$ devido à variabilidade de $\widehat{\mathbf{d}}$ em relação a $\mathbf{d}$. Por exemplo, se $\{x_1, x_2, \ldots, x_n\}$ for uma amostra aleatória de uma distribuição normal e $\widehat{\theta} = \overline{x}$, então $cv(\widehat{se}_\infty) = \frac{1}{\sqrt{2n}}$, resultando em $0,22$ para $n = 10$.

A Equação~(\ref{eq:2.19}) tem uma consequência prática importante: dados os valores típicos de $cv(\widehat{se}_\infty)$ e $\Delta$ encontrados na prática, $cv(\widehat{se}_m)$ não será muito maior que $cv(\widehat{se}_\infty)$ quando $m \geq 200$.

%-----------------------------------------------------------------------------------
\begin{table}[!htbp]
\centering
\begin{threeparttable}
\caption{O coeficiente de variação.}
\begin{tabular}{lllllll}
\hline
 &  &  m\tnote{a} $\to$ &  &  &  &  \\ 
 &  & 25 & 50 & 100 & 200 & $\infty$ \\ 
\hline
 $cv(\widehat{se}_{\infty})$\tnote{b} & 0,25 & 0,29 & 0,27 & 0,26 & 0,25 & 0,25 \\ 
\multicolumn{1}{c}{$\downarrow$} & 0,20 & 0,24 & 0,22 & 0,21 & 0,21 & 0,20 \\ 
 & 0,15 & 0,21 & 0,18 & 0,17 & 0,16 & 0,15 \\ 
 & 0,10 & 0,17 & 0,14 & 0,12 & 0,11 & 0,10 \\ 
 & 0,05 & 0,15 & 0,11 & 0,09 & 0,07 & 0,05 \\ 
 & 0,00 & 0,14 & 0,10 & 0,07 & 0,05 & 0,00 \\ 
\hline
\end{tabular}
\begin{tablenotes}[flushleft] % Garante o alinhamento exatamente abaixo do analito
    \footnotesize
    \item Fonte:~EFRON, B.; TIBSHIRANI, R. J. \textit{An Introduction to the Bootstrap}, USA: Chapman \& Hall/CRC, 1993, p.53.
	\item[a] Número de amostras \textit{bootstrap}.
	\item[b] Coeficiente de variação.
\end{tablenotes}
\label{tab:coefvaria}
\end{threeparttable}
\end{table}
%-----------------------------------------------------------------------------------

A Tabela~\ref{tab:coefvaria} compara $cv(\widehat{se}_m)$ com $cv(\widehat{se}_{\infty})$ para várias escolhas de $m$, assumindo $\Delta = 0$. Para $cv(\widehat{se}_\infty)\geq 0{,}10$, com $m=100$ o $cv(\widehat{se}_m)$ excede o ideal em no máximo \num{0,02} (Tabela~\ref{tab:coefvaria}). Efron~\cite{EFRON:1993} sugere duas regras práticas através da sua experiência:

\begin{enumerate}[label=\alph*)]
    \item Mesmo um número pequeno de replicações \textit{bootstrap}, como $m = 25$, é geralmente informativo; $m = 50$ é frequentemente suficiente para uma boa estimativa de $se_{\mathbf{d}}(\widehat{\theta})$.
    \item Segundo Efron e Tibshirani~\cite{EFRON:1993}, raramente são necessárias mais de 200 replicações para o erro padrão; para intervalos de confiança.
\end{enumerate}

A Tabela~\ref{tab:coefvaria} mostra a variação de $\widehat{se}_m$ em função do coeficiente de variação da estimativa ideal $\widehat{se}_{\infty}$ e do número de amostras \textit{bootstrap} $m$, obtidos da Equação~(\ref{eq:2.19}) assumindo $\Delta = 0$.
%-----------------------------------------------------------------------------------

\section{Reamostragem Aleatória}

O processo de reamostragem é realizado a partir da seleção aleatória de números inteiros $\{j_1, j_2, \ldots, j_n\}$, cada um com probabilidade $1/n$. Uma amostra aleatória de tamanho $n$ é definida pela seleção aleatória de $n$ unidades $\{u_1, u_2, \ldots, u_n\}$ extraídas da população \gls{u} (ver Figura~\ref{fig:amoale}), onde \gls{u} é a população de números inteiros selecionados para gerar variáveis aleatórias. 

A população \gls{v} é formada pelos valores selecionados de \gls{u} de acordo com os valores de $j$, sendo constituída por $\{v_1, v_2, \ldots, v_n\}$. Esta população é criada por um dispositivo (como um software) responsável por gerar as amostras aleatórias. Por meio dos valores $\{j_1, j_2, \ldots, j_n\}$ selecionados, pode-se determinar quais membros de \gls{v} serão escolhidos para a amostra aleatória:
\begin{equation*}
    u_1 = v_{j_1}, \quad u_2 = v_{j_2}, \quad \ldots, \quad u_n = v_{j_n}.
\end{equation*}

Sendo \gls{x_i} uma medida da unidade $u_i$ para $i = 1, 2, \ldots, n$, o conjunto de dados observados será
\begin{equation}
    \gls{x} = \{x_1, x_2, \ldots, x_n\}.
\end{equation}
%----------------------------------------------------------------------
\begin{figure}[H]%ht!
\centering
\caption{Processo de amostragem aleatória.}
\resizebox{0.25\textwidth}{!}{
\begin{tikzpicture}[
    box/.style={rectangle, draw, minimum height=1.5cm, text width=4cm, align=center},
    arrow_style/.style={-Stealth, thick},
    label_style/.style={font=\small, align=center}
]
% Nodos principais
\node[box] (populacao_v) at (0, 0) {\textbf{População $\mathbf{v}$}\\ $v_1, v_2, \ldots, v_n$};
\node[box, below=1.5cm of populacao_v] (amostra_u) {\textbf{Amostra $\mathbf{u}$}\\ $u_1, u_2, \ldots, u_n$};
\node[box, below=1.5cm of amostra_u] (medidas) {\textbf{Medidas $\mathbf{x}$}\\ $x_1, x_2, \ldots, x_n$};

% Setas e rótulos
\draw[arrow_style] (populacao_v.south) -- node[right, label_style] {Seleção por índices\\ $j_1, j_2, \ldots, j_n$\\ $u_i = v_{j_i}$} (amostra_u.north);
\draw[arrow_style] (amostra_u.south) -- node[right, label_style] {Observação\\ $x_i = $ medida de $u_i$} (medidas.north);

\end{tikzpicture}
}
\legend{\footnotesize Fonte:~\oautor}
\label{fig:amoale}
\end{figure}
%----------------------------------------------------------------------
Se as unidades da amostra $\{u_1, u_2, \ldots, u_n\}$ são observadas, obtêm-se as medidas correspondentes $\{x_1, x_2, \ldots, x_n\}$, que formam o vetor de dados \gls{x}.
%----------------------------------------------------------------------

%----------------------------------------------------------------------
\section{Teoria da Probabilidade}
\label{sec:teopro}
%----------------------------------------------------------------------
A partir de uma amostra aleatória \gls{x_al}, deve-se inferir a distribuição de probabilidade da população completa $\mathbf{X} = [X_1, X_2, \ldots, X_n]$, onde $X_i$ representa a variável aleatória associada à $i$-ésima observação. A teoria da probabilidade segue o raciocínio oposto: a partir do conhecimento da distribuição de probabilidade da população $\mathbf{X}$, deduzem-se as propriedades das amostras aleatórias e das estatísticas calculadas a partir delas.

Seja \gls{x_al2} o resultado da rolagem de um dado não viciado. Por exemplo, no lançamento de um dado cúbico de 6 faces, onde todas as faces possuem probabilidades iguais de ocorrência, obtém-se como resultado 1, 2, 3, 4, 5 ou 6. Em notação probabilística, tem-se
\begin{equation}
    \Pr\{x_{al} = k\} = \frac{1}{6}, \quad \text{para } k = 1, 2, 3, 4, 5, 6.
\end{equation}

%------------------------------------------------------------------------------------------
A quantidade aleatória que possua a probabilidade de ser selecionada como \gls{x_al2} é chamada de \textit{variável aleatória}. Na teoria da probabilidade, a notação utilizada para definir a frequência \gls{f_k} é dada por
\begin{equation}
    f_k = \Pr\{x_{al} = k\}.
\end{equation}

Na rolagem do dado não viciado, tem-se $f_k = \frac{1}{6}$ para $k = 1, 2, \ldots, 6$. A distribuição de probabilidade para a variável aleatória $x_{al}$ será denotada por \gls{f}. Aqui, \gls{f} é o vetor de frequências
\begin{equation}
    \mathbf{f} = (f_1, f_2, \ldots, f_6) = \left(\frac{1}{6}, \frac{1}{6}, \ldots, \frac{1}{6}\right).
\end{equation}

Podem-se encontrar também amostras com variáveis aleatórias \gls{x_al2} na forma da distribuição binomial de tamanho $n$ e probabilidade de sucesso $p$, denotada
\begin{equation}
    x_{al} \sim \gls{BI(n,p)}.
\end{equation}

A frequência é representada por
\begin{equation}
    f_k = \binom{n}{k} p^k (1-p)^{n-k}, \quad \text{para } k = 0, 1, 2, \ldots, n,
\end{equation}
sendo $n$ um inteiro positivo, $p$ um número entre $0$ e $1$, e $\binom{n}{k}$ o coeficiente binomial definido por
\begin{equation}
    \binom{n}{k} = \frac{n!}{k!(n - k)!}.
\end{equation}

A probabilidade pode ser definida também como no exemplo da distribuição normal (ou gaussiana). Uma variável aleatória de valor real $x_{al}$, tendo a distribuição normal com média \gls{mu} e variância \gls{sigma^2}, é descrita por
\begin{equation}
    x_{al} \sim \gls{N(mu,sigma)}.
\end{equation}

A probabilidade de $x_{al}$ pertencer a um conjunto $A \subseteq \mathbb{R}$ é dada por
\begin{equation}
    \Pr\{x_{al} \in A\} = \int_{A} \frac{1}{\sqrt{2\pi\sigma^2}} \exp\left(-\frac{1}{2}\left(\frac{x-\mu}{\sigma}\right)^2\right)dx.
\label{eq:probabili}
\end{equation}

Existem outros problemas com dimensões maiores envolvendo integrais similares à Equação~(\ref{eq:probabili}), mas não será necessário utilizar estas distribuições contínuas para o \textit{bootstrap} neste trabalho.

A expectativa da variável aleatória de valor real $x_{al}$, escrita por \gls{E(x)}, é o valor médio sobre os possíveis resultados de $x_{al}$. Assim, para a distribuição binomial:
\begin{equation}
    E(x_{al}) = \sum_{k=0}^{n} k \binom{n}{k} p^k (1 - p)^{n-k}, \quad \text{para } x_{al} \sim \text{BI}(n,p),
\end{equation}
e para a distribuição normal:
\begin{equation}
    E(x_{al}) = \int_{-\infty}^{\infty} x \frac{1}{\sqrt{2\pi\sigma^2}} \exp\left(-\frac{1}{2}\left(\frac{x-\mu}{\sigma}\right)^2\right)dx, \quad \text{para } x_{al} \sim N(\mu,\sigma^2).
\label{eq:probabili2}
\end{equation}

A expectativa, como a teoria da média, é muito ampla e inclui casos de variáveis aleatórias sem valor real definido.

Pode-se obter a variância de $x_{al}$, representada por \gls{sigma^2}. Tendo $\mu = \gls{E_f(x)}$ para $x_{al}$ sendo uma variável aleatória de valor real, a variância é definida como o valor esperado de $(x_{al} - \mu)^2$. Em outras palavras, $\sigma^2$ é o quadrado médio das distâncias da variável aleatória $x_{al}$ em relação à sua expectativa $\mu$~\cite{Canty2001}:
\begin{equation}
    \sigma^2 = E_{\mathbf{f}}[(x_{al} - \mu)^2].
\end{equation}

A variância de $x_{al} \sim N(\mu, \sigma^2)$ é igual a $\sigma^2$, e a variância de $x_{al} \sim \text{BI}(n, p)$ é igual a $np(1 - p)$. O desvio padrão da variável aleatória é definido pela raiz quadrada destas variâncias.

Para que duas variáveis aleatórias \gls{y_cal2} e \gls{z_cal} sejam independentes, deve-se ter
\begin{equation}
    E[\gls{g(y_cal)} \gls{h(z_cal)}] = E[g(y_{cal})] \cdot E[h(z_{cal})]
\end{equation}
para todas as funções $g$ e $h$. A independência implica que o resultado aleatório de $y_{cal}$ não afeta o resultado aleatório de $z_{cal}$, e vice-versa.

%-------------------------------------------------------------------------------------------------------
\section{Função de Distribuição Empírica}
\label{sec:fundisemp}
\index{A Técnica de \textit{Bootstrap}!Função de Distribuição Empírica}

Suponha que uma amostra aleatória de tamanho $n$, dada por $\gls{x} = \{x_1, x_2, \ldots, x_n\}$, seja obtida de uma função de distribuição de probabilidade \gls{d}. A função de distribuição empírica \gls{hat_d} é definida como a distribuição discreta que atribui probabilidade $\frac{1}{n}$ a cada valor observado $x_i$, para $i = 1, 2, \ldots, n$.

Em outras palavras, a distribuição empírica $\widehat{\mathbf{d}}$ é construída a partir do espaço amostral de $\gls{x}$, atribuindo a seguinte probabilidade empírica a um evento $A$:
\begin{equation}
    \widehat{\Pr}\{A\} = \frac{\#\{x_i \in A\}}{n},
\end{equation}
onde $\#\{x_i \in A\}$ denota o número de observações $x_i$ que pertencem ao conjunto $A$.

Em casos onde ocorrem repetições de valores na amostra, para evitar atribuir probabilidade $\frac{1}{n}$ a cada ocorrência repetida individualmente, utiliza-se $\widehat{\mathbf{d}}$ com as frequências relativas observadas. Para valores distintos $k$ no espaço amostral, define-se
\begin{equation}
    \widehat{f}_k = \frac{\#\{x_i = k\}}{n},
\end{equation}
onde $\#\{x_i = k\}$ representa o número de observações na amostra que são iguais a $k$. Assim, $\widehat{f}_k$ é a frequência relativa (ou probabilidade empírica) do valor $k$ na amostra.
%----------------------------------------------------------------------------------------------------
\section{Princípio \textit{Plug-in}}
\label{sec:priplu}
\index{A Técnica de \textit{Bootstrap}!Princípio \textit{Plug-in}}

O princípio \textit{plug-in} constitui uma abordagem fundamental para estimação de parâmetros estatísticos, permitindo estimar um parâmetro de interesse $\theta$ por meio de uma função avaliativa aplicada à distribuição empírica dos dados \cite{EFRON:1997}.

Formalmente, seja $\theta = t(F)$ um parâmetro definido como função da distribuição populacional $F$. O estimador \textit{plug-in} é obtido substituindo $F$ pela distribuição empírica $\hat{F}_n$:
\begin{equation}
    \hat{\theta} = t(\hat{F}_n),
\end{equation}
onde $n$ denota o tamanho da amostra.

Para o caso particular em que $\theta$ representa a expectativa populacional $\mathbb{E}_F(X)$, o estimador \textit{plug-in} corresponde à média amostral:
\begin{equation}
    \hat{\theta} = \mathbb{E}_{\hat{F}_n}(X) = \frac{1}{n} \sum_{i=1}^{n} x_i = \bar{x}.
\end{equation}

É importante destacar que o método \textit{bootstrap} constitui uma aplicação direta do princípio \textit{plug-in}, podendo ser adotado como método de estimação para diversos modelos estatísticos, incluindo modelos de regressão \cite{EFRON:1997,FERNANDEZ:1999}.

%------------------------------------------------------
\section{Técnicas Clássicas de Reamostragem}
\label{sec:teclarea}
\index{Técnicas Clássicas de Reamostragem}
%------------------------------------------------------

%Nesta seção, revisam-se duas técnicas fundamentais de reamostragem: \textit{bagging} e \textit{subagging}. Ambas utilizam amostras \textit{bootstrap}, conforme discutido na Seção~\ref{cha:tecboot}, baseando-se na distribuição empírica das probabilidades.

Nesta seção, revisam-se duas técnicas fundamentais de reamostragem: \textit{bagging} e \textit{subagging}. Ambas utilizam amostras \textit{bootstrap}, como detalhado na Seção~\ref{cha:tecboot}, baseando-se na distribuição empírica das probabilidades.

As técnicas de reamostragem são particularmente úteis em cenários com limitação amostral ou quando os dados são contaminados por ruídos instrumentais. A ideia central consiste em combinar as predições de múltiplos modelos de regressão em um único preditor agregado, processo conhecido como \textit{ensemble}~\cite{Opitz1999,GALVAO:2006}.

\begin{figure}[H]
\centering
\caption{Reamostragem com agregação de modelos de regressão.}
\label{fig:reaagremodreg}
\resizebox{0.70\textwidth}{!}{
\begin{tikzpicture}[
    box/.style={rectangle, draw, minimum height=1cm, text width=2.5cm, align=center},
    arrow_style/.style={-Stealth, thick},
    label_style/.style={font=\small, align=center}
]

% Nodo principal à esquerda
\node[box, minimum width=3cm] (conjunto_treinamento) at (0, 0) 
    {\textbf{Conjunto de}\\\textbf{Treinamento}\\ $\mathbf{X}_{\mathrm{cal}}$};

% Nodos de Subconjuntos, Regressores e Agregado
\node[box, right=2cm of conjunto_treinamento] (subconjunto1) {\textbf{Subconjunto 1}};
\node[box, right=2cm of subconjunto1] (regressor1) {\textbf{Regressor 1}};
\node[box, right=4cm of regressor1] (agregado) {\textbf{Agregado}};

\node[box, below=1cm of subconjunto1] (subconjunto2) {\textbf{Subconjunto 2}};
\node[box, below=1cm of regressor1] (regressor2) {\textbf{Regressor 2}};

\node[box, below=2cm of subconjunto1] (subconjunto_dots) {$\vdots$};
\node[box, below=2cm of regressor1] (regressor_dots) {$\vdots$};

\node[box, below=3cm of subconjunto1] (subconjunto_n) {\textbf{Subconjunto $B$}};
\node[box, below=3cm of regressor1] (regressor_n) {\textbf{Regressor $B$}};

% Setas
\draw[arrow_style] (conjunto_treinamento) -- (subconjunto1);
\draw[arrow_style] (conjunto_treinamento) -- (subconjunto2);
\draw[arrow_style] (conjunto_treinamento) -- (subconjunto_dots);
\draw[arrow_style] (conjunto_treinamento) -- (subconjunto_n);

\draw[arrow_style] (subconjunto1) -- (regressor1);
\draw[arrow_style] (subconjunto2) -- (regressor2);
\draw[arrow_style] (subconjunto_dots) -- (regressor_dots);
\draw[arrow_style] (subconjunto_n) -- (regressor_n);

\draw[arrow_style] (regressor1) -- (agregado);
\draw[arrow_style] (regressor2) -- (agregado);
\draw[arrow_style] (regressor_dots) -- (agregado);
\draw[arrow_style] (regressor_n) -- (agregado);

% Símbolos no Agregado
\node[above=0.1cm of agregado] {$\hat{\mathbf{y}}$};
\node[below=0.1cm of agregado] {$f_{\text{agg}}$};

\end{tikzpicture}
}
\legend{\footnotesize Fonte:~\oautor}
\end{figure}

A Figura~\ref{fig:reaagremodreg} ilustra o fluxo de reamostragem e agregação: o conjunto de calibração $\mathbf{X}_{\mathrm{cal}}$ é reamostrado em $B$ subconjuntos, cada um utilizado para calibrar um regressor independente. As predições individuais são subsequentemente combinadas no estágio de agregação mediante uma função $f_{\text{agg}}$.

Técnicas \textit{ensemble} integram múltiplos modelos para aprimorar o desempenho preditivo final~\cite{Opitz1999}. Neste trabalho, emprega-se a média aritmética como função de agregação, com todos os regressores utilizando Regressão Linear Múltipla (RLM), caracterizando um \textit{ensemble} homogêneo~\cite{liu2025,iterative_mishra_2023,styawan2024}.

A distinção fundamental entre \textit{bagging} e \textit{subagging} reside no esquema de amostragem: com e sem reposição, respectivamente. As Seções~\ref{sec:bagging} e~\ref{sec:subagging} detalham estas técnicas.

%----------------------------------------------------------------------------------------------------------
\subsection{\textit{Bagging}}
\label{sec:bagging}
\index{Técnicas Clássicas de Reamostragem!\textit{Bagging}} 
%----------------------------------------------------------------------------------------------------------

A técnica \textit{bagging} (\textit{Bootstrap Aggregating}) foi proposta por Breiman~\cite{Breiman1996} com o objetivo de gerar múltiplos conjuntos de calibração por meio de reamostragem \textit{bootstrap}, combinando posteriormente os resultados obtidos. Esta abordagem visa reduzir a variância do preditor final através da agregação~\cite{Breiman1996,Shabbir2024,HONG2020,nazhao2018,rathnayake2022}.

%Breiman demonstrou que o \textit{bagging} proporciona ganhos significativos de desempenho em problemas de regressão, especialmente quando o modelo é instável -- isto é, quando pequenas alterações no conjunto de calibração resultam em mudanças substanciais nas predições. Na decomposição do erro de predição em regressão, o \textit{bagging} atua principalmente na redução da componente de variância do preditor.

Breiman demonstrou que o \textit{bagging} proporciona ganhos expressivos de 
desempenho em problemas de regressão, especialmente quando o modelo é instável 
-- isto é, quando pequenas alterações no conjunto de calibração resultam em 
mudanças drásticas nas predições. Na decomposição do erro de predição, o 
\textit{bagging} atua principalmente na redução da componente de variância do 
preditor.

A técnica também se mostra eficaz em regressões multivariadas, particularmente em problemas de seleção de amostras de calibração em modelos de Regressão Linear Múltipla (RLM)~\cite{kumari2024,HARDLE:2007}. Variações na seleção de variáveis, decorrentes de flutuações nos conjuntos de calibração, são suavizadas pelo cálculo da média dos modelos resultantes~\cite{MARTINS:2006}.

\begin{center}
\begin{minipage}{0.95\textwidth}
\centering % Adiciona centralização dentro do minipage
\resizebox{0.80\textwidth}{!}{%
\begin{algorithm2e}[H]
 \dontprintsemicolon
 \linesnumbered
 \SetLine
 \Entrada{Amostra original $\mathbf{x} = (x_1, x_2, \ldots, x_n)$ com $n$ observações, estatística de interesse $\theta = s(\mathbf{x})$, número de réplicas \textit{bootstrap} $B$}
 \Saida{Estimativa \textit{bagging} $\hat{\theta}_{\text{bag}}^{*}$, distribuição das estimativas \textit{bootstrap} $\{\hat{\theta}^{*}(1), \ldots, \hat{\theta}^{*}(B)\}$}
\Inicio{
Definir o número de réplicas \textit{bootstrap} $B$\;
\Para{$b = 1$ \textbf{até} $B$}{
    Gerar amostra \textit{bootstrap} $\mathbf{x}^{*(b)}$ de tamanho $n$ com reposição a partir de $\mathbf{x}$\;
    Calcular a estatística para a réplica $b$:\;
    $\hat{\theta}^{*}(b) = s(\mathbf{x}^{*(b)})$\;
}
Agregar as réplicas via média aritmética:\;
$$\hat{\theta}_{\text{bag}}^{*} = \frac{1}{B}\sum_{b=1}^{B} \hat{\theta}^{*}(b)$$
}
\caption{Estimativa \textit{bagging}.}
\label{alg:bagging}
\end{algorithm2e}
}
\legend{\footnotesize Fonte:~\oautor}
\end{minipage}
\end{center}

%\medskip

O Algoritmo~\ref{alg:bagging} descreve o procedimento para aplicação do \textit{bagging}. Inicialmente, define-se o número $B$ de réplicas \textit{bootstrap} a serem geradas. Para cada réplica $b = 1, \ldots, B$, extrai-se uma amostra $\mathbf{x}^{*(b)}$ de tamanho $n$ com reposição da amostra original, aplica-se a função $s(\cdot)$ para obter $\hat{\theta}^{*}(b)$, e finalmente agrega-se os resultados através da média aritmética~\cite{BAUER:1999}. 

No \textit{bagging} tradicional proposto por Breiman~\cite{BREIMAN:2001}, cada amostra \textit{bootstrap} possui o mesmo tamanho da amostra original. Neste trabalho, a implementação do \textit{bagging} utilizando Regressão Linear Múltipla (RLM) é denominada \textbf{RLM-\textit{bagging}}.

%----------------------------------------------------------------------------------------------------------
\subsection{\textit{Subagging}}
\label{sec:subagging}
\index{Técnicas Clássicas de Reamostragem!\textit{Subagging}}
%----------------------------------------------------------------------------------------------------------

A técnica \textit{subagging} (\textit{subsample aggregating}) baseia-se na reamostragem sem reposição, diferindo do \textit{bagging} tradicional por não permitir repetições de observações dentro das amostras~\cite{GALVAO:2006,Buhlmann2003,revelas2024,qin2023,zhaoY2018,shi2021,zhang2017,Tao2021,paleologo2020,wei2021}. Cada modelo é construído de maneira similar ao \textit{bagging}, porém utilizando subamostras de tamanho reduzido $k < n$.
%-------------------

%-------------------
\begin{center}
\begin{minipage}{0.85\textwidth}
\centering % Adiciona centralização dentro do minipage
\resizebox{0.80\textwidth}{!}{%
\begin{algorithm2e}[H]
 \dontprintsemicolon
 \linesnumbered
 \SetLine
 \Entrada{Amostra original $\mathbf{x} = (x_1, x_2, \ldots, x_n)$ com $n$ observações, estatística de interesse $\theta = s(\mathbf{x})$, número de réplicas $B$, tamanho das subamostras $k$ (onde $k \leq n$)}
 \Saida{Estimativa subagging $\hat{\theta}_{\text{subag}}^{*}$, conjunto de estimativas $\{\hat{\theta}^{*}(1), \ldots, \hat{\theta}^{*}(B)\}$}
\Inicio{
Definir o tamanho das subamostras $k$ (tipicamente $k < n$)\;
Definir o número de réplicas $B$\;
\Para{$b = 1$ \textbf{até} $B$}{
    Gerar subamostra $\mathbf{x}^{*(b)}$ de tamanho $k$ \textbf{sem reposição} a partir de $\mathbf{x}$\;
    Calcular a estatística para a subamostra:\;
    $\hat{\theta}^{*}(b) = s(\mathbf{x}^{*(b)})$\;
}
Agregar as réplicas via média aritmética:\;
$$\hat{\theta}_{\text{subag}}^{*} = \frac{1}{B}\sum_{b=1}^{B} \hat{\theta}^{*}(b)$$
}
\caption{Estimativa \textit{subagging}.}
\label{alg:subagging}
\end{algorithm2e}
}
\legend{\footnotesize Fonte:~\oautor}
\end{minipage}
\end{center}

\medskip

O Algoritmo~\ref{alg:subagging} descreve o procedimento para aplicação do \textit{subagging}. Inicialmente, definem-se os parâmetros $k$ (tamanho das subamostras) e $B$ (número de réplicas). Para cada iteração $b = 1, \ldots, B$, gera-se uma subamostra $\mathbf{x}^{*(b)}$ de tamanho $k$ \textbf{sem reposição} a partir dos dados originais, aplica-se a função $s(\cdot)$, e finalmente agrega-se os resultados através da média aritmética.

A ausência de reposição implica que cada observação aparece no máximo uma vez por subamostra, diferenciando fundamentalmente esta técnica do \textit{bagging}. Esta característica torna o \textit{subagging} particularmente sensível à escolha do tamanho das subamostras $k$ e do número de réplicas $B$. A relação ótima entre estes parâmetros, bem como uma análise comparativa entre \textit{bagging} e \textit{subagging}, são discutidas na Seção~\ref{sec:bagxsuagg}.

%----------------------------------------------------------------------
Bühlmann~\cite{Buhlmann2003} introduziu o subagging como uma otimização do \textit{bagging} para
subamostras de dimensão $m < n$. A generalização para regressão de alta dimensionalidade foi 
formalizada por Ing e Lai~\cite{Ing2016}. Análises teóricas recentes~\cite{shi2021,revelas2024}
determinam condições necessárias e suficientes para a convergência e eficiência assintótica\footnote{Assintótica refere-se ao comportamento de algo quando uma variável tende ao infinito (ou a um limite extremo).} 
desses métodos.
%-----------------------------------------------------------------------

Neste trabalho, a implementação do \textit{subagging} utilizando Regressão Linear Múltipla (RLM) é denominada \textbf{RLM-\textit{subagging}}.

%----------------------------------------------------------------------------------------------------------
\subsection{Comparação entre \textit{Bagging} e \textit{Subagging}}
\label{sec:bagxsuagg}
\index{Técnicas Clássicas de Reamostragem!Comparação}
%----------------------------------------------------------------------------------------------------------

Embora \textit{bagging} e \textit{subagging} compartilhem o princípio de agregação de múltiplos modelos, diferem fundamentalmente no mecanismo de reamostragem empregado. O Quadro~\ref{tab:bag-vs-subag} sintetiza as principais distinções entre estas técnicas, destacando suas características operacionais e implicações para a variância dos estimadores resultantes.
%---------------------------------------------------

\begin{quadro}[ht]
\centering
\caption{Características comparativas entre \textit{bagging} e \textit{subagging}.}
\label{tab:bag-vs-subag}
\begin{tabular}{|l|c|c|}
\hline % Substitua por \toprule se usar o pacote booktabs
\textbf{Característica} & \textbf{\textit{Bagging}} & \textbf{\textit{Subagging}} \\
\hline % Substitua por \midrule se usar o pacote booktabs
Tipo de amostragem & Com reposição & Sem reposição \\
Tamanho das réplicas & $n$ & $k \leq n$ \\
Repetição de observações & Permitida & Não permitida \\
Diversidade entre réplicas & Moderada & Alta \\
Variância do estimador & Redução moderada & Redução potencialmente maior \\
Sensibilidade a parâmetros & Baixa & Alta ($k$ e $B$) \\
\hline % Substitua por \bottomrule se usar o pacote booktabs
\end{tabular}
\par\vspace{3pt}
{\footnotesize Fonte: \oautor}
\end{quadro}


Como mostra o Quadro~\ref{tab:bag-vs-subag}, a ausência de reposição no \textit{subagging} resulta em maior diversidade entre as réplicas, uma vez que cada subamostra contém um conjunto distinto de observações. Esta característica pode proporcionar maior redução de variância, mas torna a técnica mais sensível à escolha adequada dos parâmetros $k$ (tamanho das subamostras) e $B$ (número de réplicas).


A escolha entre \textit{bagging} e \textit{subagging} depende fundamentalmente das características do problema em questão. O \textit{bagging} é mais robusto e requer menos ajuste de parâmetros, sendo preferível quando se dispõe de amostras de tamanho moderado. Por outro lado, o \textit{subagging} pode ser vantajoso em cenários com amostras grandes ou quando há interesse em controlar explicitamente o tamanho das subamostras para reduzir o custo computacional~\cite{Buhlmann2003}.

Nas aplicações com Regressão Linear Múltipla (RLM) desenvolvidas neste trabalho, ambas as técnicas serão avaliadas empiricamente, permitindo uma análise comparativa de seus desempenhos preditivos em diferentes configurações amostrais.

%------------------------------------------------------
\section{Fundamentos de Engenharia de Migração e Qualidade de Software}
\label{sec:fundengsw}
\index{Revisão da literatura!Engenharia de Migração e Qualidade de Software}
%------------------------------------------------------

Assim como a calibração multivariada, a modernização de sistemas legados possui
fundamentação teórica própria. Sommerville~\cite{sommerville2011} estrutura a
engenharia de requisitos em quatro fases interdependentes --- elicitação, análise,
especificação e validação ---, oferecendo o arcabouço metodológico adotado neste
trabalho para decompor as responsabilidades do sistema Java legado em módulos
independentes, evitando que a reescrita repetisse o acoplamento entre lógica de
negócio, persistência e interface presente na implementação original.

A estratégia de migração adotada --- reescrita caixa preta --- extrai os
requisitos do comportamento observável de um sistema legado, sem reaproveitar seu
código-fonte~\cite{becker2020}. Diferentemente da reescrita caixa-branca, que
preserva estruturas internas do sistema original, essa abordagem trata o legado
exclusivamente como especificação funcional, permitindo à equipe redesenhar a
arquitetura livremente, desde que a equivalência comportamental da lógica
matématica seja formalmente demonstrada em cada etapa da migração.

A verificação dessa equivalência comportamental exige um critério de avaliação independente da implementação~\cite{pressman2016}. Testes de caixa-preta comparam entradas e saídas entre o sistema original e o sistema reescrito sem inspecionar o
código-fonte legado, estratégia adequada quando o objetivo --- como neste trabalho
--- é substituir integralmente a plataforma de execução preservando o comportamento
matemático. Essa fundamentação sustenta a suíte de 92 testes automatizados que
compara Java e Python, apresentada adiante.

%------------------------------------------------------
\subsection{Teste de Software}
\label{sec:testesoftware}
\index{Revisão da literatura!Teste de Software}
%------------------------------------------------------

O teste de software verifica se um sistema se comporta conforme especificado,
revelando defeitos antes da operação~\cite{pressman2016}. Neste trabalho,
empregam-se o teste unitário, que exercita funções isoladas; o teste de
caixa-preta, que confronta entradas e saídas sem inspecionar o código
interno~\cite{pressman2016,jorgensen2014}; e o teste de carga e
\textit{stress}, que mede tempo e memória sob grandes volumes de dados. A
automação com \textit{frameworks} como o \texttt{pytest}~\cite{pytest2024}
permite repetir essas verificações a cada alteração. Na migração, o teste de
caixa-preta serve para verificar a equivalência entre o sistema legado e o
reescrito, pois compara apenas o comportamento observável.